%%% AAMAS 2027 submission (Research Paper Track), anonymized.
\documentclass[sigconf,anonymous]{aamas}

\usepackage{balance}
\usepackage{xspace}
\usepackage{algorithm}
\usepackage{algorithmic}

%%% AAMAS-2027 copyright block (do not change!)
\makeatletter
\gdef\@copyrightpermission{
  \begin{minipage}{0.2\columnwidth}
   \href{https://creativecommons.org/licenses/by/4.0/}{\includegraphics[width=0.90\textwidth]{by}}
  \end{minipage}\hfill
  \begin{minipage}{0.8\columnwidth}
   \href{https://creativecommons.org/licenses/by/4.0/}{This work is licensed under a Creative Commons Attribution International 4.0 License.}
  \end{minipage}
  \vspace{5pt}
}
\makeatother

\setcopyright{ifaamas}
\acmConference[AAMAS 2027]{Proc.\@ of the 26th International Conference
on Autonomous Agents and Multiagent Systems (AAMAS 2027)}{3--7 May 2027}{Hanoi, Vietnam}{M.~Baldoni, F.~Fang, W.~Yeoh, N.~Yorke-Smith (eds.)}
\copyrightyear{2027}
\acmYear{2027}
\acmDOI{}
\acmISBN{}

\submissionType{Research Paper Track}

% shared macros (main paper)
\newcommand{\ind}{\mathbf 1}
% pointers into the technical appendix (supplementary material, lettered sections)
\newcommand{\appTwo}{Appendix~A\xspace}
\newcommand{\appFamily}{Appendix~B\xspace}
\newcommand{\appCoverage}{Appendix~C\xspace}
\newcommand{\appFifteen}{Appendix~D\xspace}
\newcommand{\appAllsizes}{Appendix~E\xspace}
\newcommand{\appComplexity}{Appendix~F\xspace}
\newcommand{\appDistance}{Appendix~G\xspace}
\newcommand{\appDefect}{Appendix~H\xspace}


\title[Subjective SD-EF1 and Market EF1]{Subjective SD-EF1 and Market EF1:\\Two-Agent Existence and Barriers for Three Agents}

\author{Anonymous Authors}
\affiliation{\institution{Anonymous}\city{Anonymous}\country{Anonymous}}
\email{anonymous@example.com}


\begin{CCSXML}
<ccs2012>
<concept>
<concept_id>10003752.10010070.10010099.10010100</concept_id>
<concept_desc>Theory of computation~Algorithmic game theory</concept_desc>
<concept_significance>500</concept_significance>
</concept>
<concept>
<concept_id>10003752.10003777.10003779</concept_id>
<concept_desc>Theory of computation~Problems, reductions and completeness</concept_desc>
<concept_significance>300</concept_significance>
</concept>
<concept>
<concept_id>10010147.10010178.10010219.10010220</concept_id>
<concept_desc>Computing methodologies~Multi-agent systems</concept_desc>
<concept_significance>300</concept_significance>
</concept>
</ccs2012>
\end{CCSXML}

\ccsdesc[500]{Theory of computation~Algorithmic game theory}
\ccsdesc[300]{Theory of computation~Problems, reductions and completeness}
\ccsdesc[300]{Computing methodologies~Multi-agent systems}

\begin{abstract}
Barman et al.\ asked whether indivisible goods can always be allocated so that every agent is stochastic-dominance envy-free up to one good (SD-EF1) with respect to its own ranking, while the allocation is EF1 for a common market valuation.
We settle the two-agent case in a strong form. From the two rankings alone we construct $\lceil m/2\rceil+1$ subjectively SD-EF1 allocations such that every inclusion-monotone market valuation makes one of them market EF1, and $O(\log m)$ value queries find it; the result extends to chores and to mixed items with a shared classification. Every such family needs $\Omega(\sqrt m)$ allocations, and on all sampled hard instances with $m\le 8$ exact computations show that $\lceil m/2\rceil$ are needed, so our family is optimal there up to its redundant mirror-image endpoint.
For three agents whose market order coincides with one agent's ranking, we prove existence for up to fifteen goods by an exhaustive computation, and for every number of goods under a structural condition on the remaining ranking.
We then explain why the natural matching-based route to three agents stalls: completing a prescribed bundle is NP-complete, repairing one can require a linear number of replacements, requiring balance in the owner's own ranking fails from thirteen goods on, and pairs of allocations with disjoint market defects need not exist, even for two agents.
General three-agent existence remains open; exact computations on thousands of hard profiles find no counterexample.
\end{abstract}

\keywords{Fair division; Envy-freeness up to one good; Stochastic dominance; Market values; Computational complexity}

\begin{document}
\maketitle

%%%%%%%%%%%%%%%%%%%%%%%%%%%%%%%%%%%%%%%%%%%%%%%%%%%%%%%%%%%%%%%%%%%%%%%%
\section{Introduction}
Fair division problems often carry two kinds of value information. Agents have \emph{subjective} preferences over the goods, and the market attaches a public value to each good: an appraisal in an inheritance, a list price, an expert's estimate. Envy-freeness up to one good (EF1)~\citep{lipton2004ef1,budish2011} is the standard fairness notion for indivisible goods, but against which valuation should it be measured? \citet{barman2025market} propose to ask for both. Near-equal division under market values can be a legal or normative requirement, and the agents' reports can be read as noisy estimates of the market values; an allocation that is fair for both is robust to disputes about which valuation is the right one.

Formally, each agent~$i$ has a strict ranking $\sigma_i$ of the goods, standing for an additive utility, and all agents share a nonnegative additive market valuation~$v$. Agent~$i$ is \emph{SD-EF1} towards~$j$ if it is EF1 for every additive utility consistent with~$\sigma_i$; equivalently, its bundle stochastically dominates $j$'s bundle after one good is deleted~\citep{aziz2015ordinal}. \citet{barman2025market} show that SD-EF1 on both sides is impossible already for two agents and seven goods, whereas EF1 for the agents together with SD-EF1 for the market can always be achieved. Their Open Question~1 asks about the remaining direction:
\begin{quote}\emph{Is there always an allocation that is SD-EF1 for the agents' rankings and EF1 for the market valuation?}\end{quote}
The authors remark that this ``seems to require considerably novel ideas''.

\paragraph{Contributions.}
\begin{enumerate}
\item \emph{Two agents (Section~\ref{sec:two}).} The answer is yes, in a strong form. From the two rankings alone we build a sequence of $\lceil m/2\rceil+1$ subjectively SD-EF1 allocations, consecutive ones differing by a swap (a single-good transfer if $m$ is odd), such that \emph{every inclusion-monotone} market valuation makes one of them market EF1; $O(\log m)$ value queries find it. The same path handles chores, and mixed items when agents and market agree on which items are goods. For monotone markets a family of size $\Omega(\sqrt m)$ is unavoidable, and exact computation shows that ours, with its redundant mirror-image endpoint removed, is a minimum family on all sampled hard instances with $m\le 8$.
\item \emph{Three agents, one ranking equal to the market order (Section~\ref{sec:three}).} Existence for up to fifteen goods by an exhaustive computer-assisted argument, and for every number of goods when the third ranking is ``one-singleton'', a class containing all rankings obtained from the market order by disjoint adjacent transpositions.
\item \emph{Barriers (Section~\ref{sec:barriers}).} A natural route to three agents fixes one agent's bundle as a perfect matching between market triples and ranking triples and completes the rest. We show that this route stalls: completing a prescribed matching is NP-complete; repairing it can require $\Omega(m)$ replacements; requiring the bundle to be balanced in its owner's own ranking fails from thirteen goods on; and pairs of allocations whose market defects are disjoint, which suffice for additive markets, need not exist, even for two agents.
\end{enumerate}
General three-agent existence remains open. Section~\ref{sec:discussion} reports exact computations that find no counterexample among thousands of hard profiles, and lists open questions. Table~\ref{tab:scope} states the scope of each result; the results concern different restrictions and should not be conflated.

\begin{table*}[t]
\caption{Scope of the results. ``Proof'' means a written proof; ``computer-assisted'' a proof with an exhaustive finite step; ``evidence'' is not a theorem.}
\label{tab:scope}
\small
\begin{tabular}{@{}p{.33\textwidth}p{.46\textwidth}p{.17\textwidth}@{}}
\toprule
Setting & Guarantee or boundary & Nature\\
\midrule
Two agents, any $m$, any monotone market & SD-EF1 and market EF1; one ranking-based family of $\lceil m/2\rceil+1$ allocations, $O(\log m)$ queries & Proof (Thm.~\ref{thm:two})\\
Two agents, monotone markets & every cover has $\Omega(\sqrt m)$ members; minimum $\lceil m/2\rceil$ on sampled hard profiles, $m\le8$ & Proof; computation\\
Three agents, $\sigma_0=\pi$, $m\le15$ & market-block-balanced SD-EF1 allocation exists; own-bundle balance exactly up to $m=12$ & Computer-assisted (Thm.~\ref{thm:twelve})\\
Three agents, $\sigma_0=\pi$, $\sigma_2$ one-singleton, any $m$ & existence in $O(m)$ time & Proof (Thm.~\ref{thm:onesingleton})\\
Prescribed reset/suffix matching & completion is NP-complete & Proof (Thm.~\ref{thm:hard})\\
SAT-reduction profiles, matching free & linear-time construction, at most one replacement & Proof (Thm.~\ref{thm:repair})\\
Connected family, own-bundle balance & $\Omega(m)$ replacements may be needed & Proof (Thm.~\ref{thm:distance})\\
Two allocations with disjoint market defects & suffice for additive markets; need not exist ($m=12$; two agents, $m=11$) & Proof; enumeration\\
General three-agent Open Question~1 & open; no counterexample among 4{,}087 + 102 hard profiles & Evidence\\
\bottomrule
\end{tabular}
\end{table*}

\paragraph{Why two agents are different.}
For two agents one can walk from an allocation to its mirror image; the market envy must change sign somewhere along the walk, and the point is to keep both agents' prefix conditions intact throughout (Lemmas~\ref{lem:jump}--\ref{lem:flip}). For three agents no such one-dimensional path is available: on hard instances the complex of subjectively fair allocations has holes, so the topological transfer principle of \citet{dupre2026balanced}, which needs a simply connected complex for three agents, does not apply directly (Section~\ref{sec:discussion}).

\paragraph{Related work.}
The simultaneous market model is due to \citet{barman2025market}. Their block characterization connects market SD-EF1 to category constraints, for which \citet{biswas2018cardinality} obtain EF1 guarantees; such cardinal guarantees do not supply the prefix inequalities required here. EF1 goes back to \citet{lipton2004ef1,budish2011}, stochastic-dominance fairness to \citet{aziz2015ordinal}; see \citet{suksompong2021constraints} for constrained fair division. Balanced allocations, the special case of uniform market values, are studied by \citet{dupre2026balanced} and \citet{cookson2026balanced}.

For two agents, our starting point is the alternating-pair partition of \citet[Lemma~4.1]{kyropoulou2020group}; the new ingredient is a flip order that preserves the two \emph{directed} prefix inequalities while crossing a market-value interval. One-dimensional intermediate-value arguments recur in two-agent fair division~\citep{kyropoulou2020group,dupre2026balanced}; see also \citet{goldberg2022halving,dupre2026necklace} for partitions that are fair for several valuations at once. \citet{bu2023allocator} prove non-constructively that two agents with monotone valuations admit an allocation that is EF1 both for the agents and for an allocator whose two evaluation functions may differ (Remark~3.6 there; Remark~4.8 gives polynomial time if one side is additive). With a single common monotone market, ordinary doubly-EF1 also follows constructively from an Exact1 partition~\citep[Appendix~A]{goldberg2022halving} for agent~1 and the market, followed by agent~2's choice. Our theorem is incomparable with these: it needs one common market, but strengthens the agents' side to SD-EF1 (one allocation must serve every additive utility consistent with each ranking) and provides one ranking-only family serving every monotone market.

For three agents the question is a one-sided prefix-balancing problem for three rankings plus a market condition. Two-sided balancing of three permutations can require logarithmic discrepancy~\citep{newman2012beck}, and group EF1 uses discrepancy theory~\citep{manurangsi2022groups}, but these bounds concern two-sided constraints. \citet{prakash2025completion} study completing frozen partial allocations and prove EF1-completion hard even for binary additive preferences; our completion problem fixes a whole bundle of one of three agents under market-block balance and strict SD-EF1, with additional reset and suffix structure.

%%%%%%%%%%%%%%%%%%%%%%%%%%%%%%%%%%%%%%%%%%%%%%%%%%%%%%%%%%%%%%%%%%%%%%%%
\section{Model and Prefix Constraints}\label{sec:model}
Let $M$ be a set of $m$ goods and $A=(A_i)_{i\in N}$ a complete allocation. Agent~$i$ has a strict ranking~$\sigma_i$ (best first); $P_i(t)$ denotes its $t$ highest ranked goods. Ties can be broken arbitrarily: satisfying all prefixes of a refinement satisfies the prefixes of the weak order.

\paragraph{Subjective fairness.} Agent~$i$ is \emph{SD-EF1} towards~$j$ if $A_j=\varnothing$ or, after deleting one good of~$A_j$, $A_i$ weakly dominates it at every prefix of~$\sigma_i$. An allocation is \emph{subjectively SD-EF1} if this holds for all ordered pairs.

\begin{lemma}[Prefix form]\label{lem:prefix}
An allocation is subjectively SD-EF1 iff $|A_j\cap P_i(t)|\le |A_i\cap P_i(t)|+1$ for all $i\ne j$ and all~$t$.
\end{lemma}
\begin{proof}
Necessity: deleting a good lowers a prefix count by at most one. Sufficiency: delete $A_j$'s $\sigma_i$-best good; before it, $A_j$ has no good in the prefix, afterwards the inequality applies.
\end{proof}

\paragraph{Market fairness.} A market valuation $v:2^M\to\mathbb R_{\ge0}$ is \emph{inclusion-monotone} if $S\subseteq T$ implies $v(S)\le v(T)$; additive valuations are the main special case. An allocation is \emph{market EF1} if for all $i\ne j$ with $A_j\ne\varnothing$ some $g\in A_j$ satisfies $v(A_i)\ge v(A_j\setminus\{g\})$. For additive $v$ deleting the most valuable good is optimal, which gives the usual $v(A_i)\ge v(A_j)-\max_{g\in A_j}v(g)$.

\paragraph{Market blocks.} Let $\pi$ be an order of the goods by nonincreasing market value. An allocation is \emph{market-block-balanced} if every full block of $n$ consecutive goods of~$\pi$ gives one good to each agent, and the final short block at most one. All market-prefix counts then differ by at most one, so the allocation is EF1 for every additive $v$ consistent with~$\pi$. For a strict common ranking, subjective SD-EF1 forces this structure~\citep{barman2025market}.

%%%%%%%%%%%%%%%%%%%%%%%%%%%%%%%%%%%%%%%%%%%%%%%%%%%%%%%%%%%%%%%%%%%%%%%%
\section{Two Agents}\label{sec:two}
In this section the agents are $1,2$ and no relation between the rankings and the market is assumed.

\begin{theorem}\label{thm:two}
For two agents with strict rankings $\sigma_1,\sigma_2$ there is a sequence $A^0,\dots,A^T$ of $T+1\le\lceil m/2\rceil+1$ subjectively SD-EF1 allocations, computed from the rankings alone and obtained one from the next by $T$ swaps (single-good transfers if $m$ is odd), with $A^T$ equal to $A^0$ with the two bundles exchanged, such that for every inclusion-monotone market valuation~$v$ some $A^t$ is market EF1. Given a value oracle for~$v$, such an allocation and the good witnessing EF1 are found with at most $3+2\lceil\log_2\lceil m/2\rceil\rceil$ value queries and $O(m\log m)$ further time.
\end{theorem}

The quantifiers matter: \emph{one} ranking-dependent family serves \emph{every} monotone valuation. For $A=(A_1,A_2)$ let $U$ be the set of allocations in which agent~2 fails market EF1 towards agent~1, and $L$ those in which agent~1 fails towards agent~2. If $A\in U$ then $v(A_2)<v(A_1\setminus\{g\})\le v(A_1)$ for any $g\in A_1$, and symmetrically $A\in L$ implies $v(A_1)<v(A_2)$; hence $U\cap L=\varnothing$, and $A$ is market EF1 iff it lies in neither.

\begin{lemma}[No jump]\label{lem:jump}
For every inclusion-monotone $v$, a single swap or transfer cannot lead from $U$ to $L$ or from $L$ to $U$.
\end{lemma}
\begin{proof}
Consider a swap $(X\cup\{g\},Y\cup\{h\})\to(X\cup\{h\},Y\cup\{g\})$. If the first allocation is in $U$ and the second in $L$, then deleting $g$ from agent~1's bundle does not satisfy agent~2 in the first, so $v(X)>v(Y\cup\{h\})$, and deleting the same $g$ from agent~2's bundle does not satisfy agent~1 in the second, so $v(Y)>v(X\cup\{h\})$. Monotonicity gives $v(X)>v(Y\cup\{h\})\ge v(Y)>v(X\cup\{h\})\ge v(X)$, a contradiction. If the first is in $L$ and the second in $U$, deleting $h$ gives $v(Y)>v(X\cup\{g\})\ge v(X)>v(Y\cup\{g\})\ge v(Y)$. For a transfer $(X\cup\{g\},Y)\to(X,Y\cup\{g\})$ from $U$ to $L$, deleting $g$ gives $v(X)>v(Y)$ and then $v(Y)>v(X)$. From $L$ to $U$ the unrepaired envies give $v(X\cup\{g\})<v(Y)$ and $v(Y\cup\{g\})<v(X)$, so $v(Y)\le v(Y\cup\{g\})<v(X)\le v(X\cup\{g\})<v(Y)$. Transfers towards agent~1 are symmetric.
\end{proof}

\paragraph{The path.} If $m$ is odd, append a bottom-ranked good of market value zero to both rankings. Pair consecutive goods of $\sigma_1$ by \emph{red} edges and of $\sigma_2$ by \emph{blue} edges. The union is a disjoint union of even alternating cycles (a pair that is both red and blue forms a cycle of length two). A proper two-colouring splits every red and blue pair, and both assignments of its colour classes are subjectively SD-EF1~\citep{kyropoulou2020group}. Write $w_i(t)=|A_1\cap P_i(t)|-|A_2\cap P_i(t)|$; by Lemma~\ref{lem:prefix} the \emph{directed} conditions are $w_1(t)\ge-1$ and $w_2(t)\le1$. Flipping the cycles one after another must keep both; Algorithm~\ref{alg:path} summarizes the construction.

\begin{lemma}[Ordered cycle flip]\label{lem:flip}
A properly coloured alternating cycle of $2\ell$ goods can be flipped in $\ell$ red-edge swaps such that both directed prefix conditions hold throughout.
\end{lemma}
\begin{proof}
A cycle of length two needs one swap. Otherwise let $B_0$ be the cycle's blue pair that comes first in $\sigma_2$, and label the cycle $c_1,\dots,c_{2\ell}$ so that $B_0=\{c_{2\ell},c_1\}$ with $c_1$ held by agent~1, and the red pairs are $\{c_{2j-1},c_{2j}\}$; the other blue pairs are $B_j=\{c_{2j},c_{2j+1}\}$. Swap the red pairs in increasing~$j$. All other cycles stay properly coloured, and after $j<\ell$ swaps every red pair is still split, so $w_1\ge-1$. Among the blue pairs of the cycle, $B_0$ now belongs entirely to agent~2 and $B_j$ entirely to agent~1, the others being split. Because $B_0$ precedes $B_j$ in $\sigma_2$, the value of $w_2$ at pair boundaries is $0$, then $-2$, then $0$; inside a pair it moves by one step, so $w_2\le1$ throughout. The last swap restores a proper colouring.
\end{proof}

\begin{proof}[Proof of Theorem~\ref{thm:two}]
Fix a proper colouring $(C,C')$ and let $A^0=(C,C')$. Flip the cycles one after another (Lemma~\ref{lem:flip}); this gives $A^0,\dots,A^T$ with $T\le\lceil m/2\rceil$, one swap per red pair, all subjectively SD-EF1, and $A^T=(C',C)$. Only the rankings are used. Fix a monotone $v$. If $v(C)\ge v(C')$ then $A^0\notin L$ and $A^T\notin U$, since membership in $L$ (resp.\ $U$) forces the corresponding strict value inequality; otherwise $A^0\notin U$ and $A^T\notin L$. If neither endpoint is market EF1, then in the first case $A^0\in U$ and $A^T\in L$, and at the first index leaving $U$ Lemma~\ref{lem:jump} forbids entering $L$, so that allocation is market EF1; the other case is symmetric. For odd $m$ the dummy is bottom-ranked with value zero, so deleting it changes neither the real prefixes nor market EF1, and a swap involving it becomes a transfer.

For the algorithm, query both bundle values at $A^0$; equal values give envy-freeness, and $A^T$ has the reversed values. Otherwise bisect on the sign of $v(A_1)-v(A_2)$, keeping indices with opposite signs (the sign need not be monotone along the path), until two adjacent states remain. Name the bundles so that the first is higher at the earlier state $(X\cup\{g\},Y\cup\{h\})$, where $h$ is absent for a transfer. One more query of $v(X)$ decides: if $v(X)\le v(Y\cup\{h\})$, the earlier state is market EF1 by deleting~$g$; otherwise $v(X\cup\{h\})\ge v(X)>v(Y\cup\{h\})\ge v(Y)$, so the later state is market EF1 by deleting~$g$ from the other bundle. Envy in the other direction is absent at both states. Recording each good's initial owner and flip time expands any $A^t$ in $O(m)$ time; no subjective values are queried.
\end{proof}

\begin{algorithm}[t]
\caption{\textsc{ReversalPath}$(\sigma_1,\sigma_2)$ (rankings only)}\label{alg:path}
\begin{algorithmic}[1]
\STATE if $m$ is odd, append a dummy good to both rankings
\STATE red pairs $\gets$ consecutive pairs of $\sigma_1$; blue pairs $\gets$ consecutive pairs of $\sigma_2$
\STATE split red$\,\cup\,$blue into alternating cycles; give the goods of each cycle to agents $1,2,1,2,\dots$ along it; output $A^0$
\FOR{each cycle $Z$ (any order)}
  \STATE $B_0\gets$ blue pair of $Z$ that comes first in $\sigma_2$; $c_1\gets$ its endpoint held by agent~1
  \STATE walk $c_1\xrightarrow{\text{red}}c_2\xrightarrow{\text{blue}}c_3\xrightarrow{\text{red}}c_4\cdots$ around $Z$
  \FOR{$j=1,\dots,|Z|/2$}
     \STATE swap the owners of the red pair $\{c_{2j-1},c_{2j}\}$; output the new allocation
  \ENDFOR
\ENDFOR
\end{algorithmic}
\end{algorithm}

\paragraph{Example.} Let $\sigma_1=(a,b,c,d,e,f)$ and $\sigma_2=(b,c,d,e,f,a)$. The red pairs are $ab,cd,ef$ and the blue pairs $bc,de,fa$; together they form one 6-cycle. The earliest blue pair is $bc$, so the red pairs are flipped in the order $cd$, $ef$, $ab$. The rows below give the owner of each good along the path (columns in the order of $\sigma_1$).
\begin{center}\small
\begin{tabular}{@{}lcccccc@{}}
\toprule
 & $a$ & $b$ & $c$ & $d$ & $e$ & $f$\\
\midrule
$A^0$ & 1 & 2 & 1 & 2 & 1 & 2\\
$A^1$ & 1 & 2 & \textbf{2} & \textbf{1} & 1 & 2\\
$A^2$ & 1 & 2 & 2 & 1 & \textbf{2} & \textbf{1}\\
$A^3$ & \textbf{2} & \textbf{1} & 2 & 1 & 2 & 1\\
\bottomrule
\end{tabular}
\end{center}
Every row satisfies $w_1\ge-1$ and $w_2\le1$. The intermediate rows are not balanced along $\sigma_2$ ($w_2$ reaches $-2$ in $A^1$ and $A^2$); the directed condition $w_2\le1$ permits this, whereas splitting every pair of both matchings would not.

\paragraph{Chores and mixed items.}
The same path handles chores. If agents rank chores by decreasing subjective cost and $c$ is an identical monotone market cost, apply Theorem~\ref{thm:two} with $c$ as a valuation and return the bundles in exchanged order: the goods inequalities $|B_j\cap P|\le|B_i\cap P|+1$ and $c(B_i)\ge c(B_j\setminus\{g\})$ become $|A_i\cap P|\le|A_j\cap P|+1$ and $c(A_i\setminus\{g\})\le c(A_j)$. For mixed items, toggling the owner of every chore turns each agent's utility difference into a difference of ``benefit bundles'', and the no-jump argument survives for a market utility that increases with goods and decreases with chores; this yields the analogous statement when agents and market share one classification into goods and chores, again with $O(\log m)$ queries (\appTwo). The shared classification cannot be dropped: with two items of subjective utilities $(2,-1)$ for both agents and $v(S)=|S|$, subjective EF1 forces the items together while market EF1 forces them apart.

\subsection{How many allocations are needed?}\label{sec:family}
Call a set of allocations a \emph{cover} of a class of market valuations if every valuation in the class makes some member market EF1. Theorem~\ref{thm:two} provides a cover of size $\lceil m/2\rceil+1$ consisting of subjectively SD-EF1 allocations; since $A^T$ and $A^0$ differ only by the names of the bundles, they behave identically for every market valuation, so $\lceil m/2\rceil$ different behaviours suffice.

\begin{proposition}\label{prop:size}
For $m\ge2$, every cover of the inclusion-monotone valuations on $m$ goods, by any set of allocations, has at least $(2^{m-1}-1)/(\binom{m}{\lfloor m/2\rfloor}-1)=\Omega(\sqrt m)$ members.
\end{proposition}
\begin{proof}
For $W\subseteq M$ with $|W|=2h\ge2$ let $v_W(S)=1$ if $|S\cap W|\ge h$ and $0$ otherwise; $v_W$ is monotone. An allocation is EF1 for $v_W$ iff it splits $W$ evenly: if both bundles hold $h$ goods of $W$ their values are equal; if, say, $|A_1\cap W|\ge h+1$, then $v_W(A_2)=0$ while deleting any one good from $A_1$ leaves at least $h$ goods of $W$, so agent~2 is not EF1. An allocation with $|A_1|=a$ splits exactly $\sum_{j\ge1}\binom aj\binom{m-a}j=\binom ma-1$ such sets (Vandermonde), at most $\binom m{\lfloor m/2\rfloor}-1$, and there are $2^{m-1}-1$ even subsets of size at least two.
\end{proof}

Thus no constant-size cover exists for monotone markets, in contrast to additive markets, where a single market-block-balanced allocation covers all valuations whenever it exists. Two computations indicate how far Theorem~\ref{thm:two} is from optimal (details in \appFamily). Call a pair of rankings \emph{hard} if no subjectively SD-EF1 allocation is market-block-balanced for the identity market order, so that no single allocation covers the additive valuations. \emph{(i)} On 20 sampled hard profiles with $m\in\{7,8\}$, an exact hitting-set computation shows that the smallest cover made of subjectively SD-EF1 allocations has exactly $\lceil m/2\rceil$ members, i.e.\ the family of Theorem~\ref{thm:two} without $A^T$ is a minimum cover. \emph{(ii)} On 12 sampled hard profiles with $m\in\{9,10,11\}$, exact linear programs confirm that the family covers all monotone valuations and that removing any member (up to exchanging the two bundles) destroys this. For additive markets, by contrast, two allocations covered every hard profile we tried ($m\le13$, hundreds of random profiles and an adversarial search); whether two always suffice is open (Section~\ref{sec:discussion}).

%%%%%%%%%%%%%%%%%%%%%%%%%%%%%%%%%%%%%%%%%%%%%%%%%%%%%%%%%%%%%%%%%%%%%%%%
\section{Three Agents with a Market-Aligned Ranking}\label{sec:three}
From now on the agents are $0,1,2$, the goods are numbered in market order $\pi=(0,\dots,m-1)$, and $\sigma_0=\pi$. Market-block balance then protects agent~0 and gives market EF1 for every additive valuation, so we look for a \emph{market-block-balanced subjectively SD-EF1} allocation. We call the assertion that one always exists the \emph{market-aligned conjecture}; it settles Open Question~1 whenever some agent's ranking equals the market order, but the restriction $\sigma_0=\pi$ is real. We write $(D,S,X)=(A_0,A_1,A_2)$. For $m=3q$, a good is an edge between its market triple and its $\sigma_1$ triple; $S$ is \emph{own-bundle balanced} if it takes one good of every $\sigma_1$ triple, and then $S$ is a perfect matching of this bipartite multigraph. Let $s,x$ count the goods of $S,X$ in a prefix~$P$ of length~$t$. Once $S$ is fixed, the remaining subjective constraints are exactly
\begin{align}
 t-2s-1\le x\le s+1 &\quad(P\text{ a }\sigma_1\text{ prefix}),\label{eq:one}\\
 x\ge\max\{s-1,\lfloor(t-s)/2\rfloor\} &\quad(P\text{ a }\sigma_2\text{ prefix}),\label{eq:two}
\end{align}
and each residual market pair sends one good to $X$ and one to $D$.

\begin{theorem}[Computer-assisted]\label{thm:twelve}
For three agents with $\sigma_0=\pi$ and $m\le15$ there is a market-block-balanced subjectively SD-EF1 allocation. For $m\le12$ it can be chosen own-bundle balanced for agent~1; for $m\le15$, after padding to a multiple of three, agent~1 holds two goods of at most one $\sigma_1$ triple.
\end{theorem}

\begin{proof}[Proof outline]
For $m=12$ relabel the goods so that $\sigma_1$ is the identity. Agent~0 accepts every rainbow allocation of the market triples in any order, so the market information reduces to an unordered partition~$Q$ of the goods into four triples: there are $12!/((3!)^44!)=15{,}400$. For each~$Q$ enumerate its $6^4$ rainbow allocations and keep the family $\mathcal F(Q)$ that satisfies agent~1's constraints and own-bundle balance. It remains to cover every ranking of agent~2 with a candidate that may depend on that ranking. After a prefix of $\sigma_2$ with goods $T$, let $L\subseteq\mathcal F(Q)$ be the candidates that survived all earlier prefix constraints of agent~2. The pair $(T,L)$ is a sufficient state, and
\[
\mathrm{cov}(T,L)=\begin{cases}\text{false} & L=\varnothing,\\ \text{true} & \exists A\in L:\ |A_2\cap T|\ge3,\\ \bigwedge_{g\notin T}\mathrm{cov}(T\cup g,L_g) & \text{else,}\end{cases}
\]
where $L_g$ keeps the members of $L$ with $|A_j\cap(T\cup g)|\le|A_2\cap(T\cup g)|+1$ for $j=0,1$. The middle case is safe because a candidate that already gives agent~2 three goods covers every continuation (all bundles have four goods); different branches may use different candidates. An exact integer-bitset implementation covers all 15{,}400 partitions (70{,}831{,}414 recursive calls, 43{,}515{,}662 memoized states, candidate families of 64 to 1{,}296 allocations); an independently written C++ checker covers all 15{,}400 partitions again and reproduces the family sizes and the number of memoized states exactly; a Python implementation agrees on 40 complete partition problems; and a negative control that also demands rainbow bundles for the other two agents on the $\sigma_1$ triples leaves $12{,}150$ partitions uncovered, as it must. Smaller instances are padded with bottom-ranked goods, which preserves all original prefixes and blocks. For $m\le15$ a second verifier adds two exact accelerations. A \emph{core} at $(T,L)$ is a subfamily $K\subseteq L$ that already covers every continuation; the union of the children's cores is a core of the parent, and a stored core may be reused at any later state $(T,L')$ with $K\subseteq L'$, because adding candidates never destroys coverage. Second, with $\sigma_1$ the identity, exchanging goods $0,1$ or permuting goods $12,13,14$ maps the candidate families of a partition onto those of its image: these permutations only alter prefixes of length $1$, $13$ and $14$; the first is always fair, and the other two are implied by the prefix of length $12$ (there agent~1 misses at most one of the last three goods). The verifier covers all $15!/((3!)^55!)=1{,}401{,}400$ partitions (82{,}670 orbits, $1.6\cdot10^{10}$ recursive calls) using own-balanced candidates when possible and one-shift candidates otherwise; a separate implementation keeping full states completes ten whole partition problems. \appCoverage gives the algorithm, its correctness argument and the counts.
\end{proof}

\paragraph{Own-bundle balance fails from thirteen goods.}
For $m=13$ take
\begin{align*}
 \sigma_1&=(0,1,3,12,2,6,9,10,4,7,5,8,11),\\
 \sigma_2&=(3,2,0,1,6,5,12,4,8,11,7,9,10).
\end{align*}
Of the $3{,}888$ market-block-balanced allocations, $206$ are subjectively SD-EF1, for example $A_0=\{0,3,7,9\}$, $A_1=\{1,4,6,10\}$, $A_2=\{2,5,8,11,12\}$.
\begin{proposition}\label{prop:thirteen}
No market-block-balanced subjectively SD-EF1 allocation of this profile is own-bundle balanced.
\end{proposition}
\begin{proof}
Relabel the goods by their position in $\sigma_1$. The market blocks become $Q_A=\{0,1,4\}$, $Q_B=\{2,8,10\}$, $Q_C=\{5,9,11\}$, $Q_D=\{6,7,12\}$, $Q_E=\{3\}$, and $\sigma_2$ becomes $(2,4,0,1,5,10,3,8,11,12,9,6,7)$. Own-bundle balance gives $|S|=4+\mathbf 1[12\in S]$ and market balance $|S|=4+\mathbf 1[3\in S]$, so $3\in S$ iff $12\in S$. If $3,12\in S$, then $Q_D$ excludes $6,7$ and the third $\sigma_1$ triple forces $8\in S$; otherwise $S$ takes $6$ or $7$, hence not $8$, and its $Q_B$-good is $2$ or $10$. This leaves three forms: (I) $S=\{a,3,8,v,12\}$, $a\in\{0,1\}$, $v\in\{9,11\}$; (II) $S=\{2,4,u,v\}$, $u\in\{6,7\}$, $v\in\{9,11\}$; (III) $S=\{a,5,u,10\}$, $a\in\{0,1\}$, $u\in\{6,7\}$. In case I the first eleven goods of $\sigma_2$ contain all of $S$ but at most three goods of $X$ (agent~2's $Q_D$-good is $6$ or $7$). In case II the first two goods of $\sigma_2$ lie in $S$. In case III the first six goods of $\sigma_2$ contain three $S$-goods and agent~2's $Q_A$-good, so $2\in X$; the first five goods of $\sigma_1$ then hold one $S$-good and two $X$-goods, so $3\in D$ and $|X|=4$; the first eleven goods of $\sigma_2$ contain three $S$-goods, so they contain all of $X$ and $12\in X$; the first nine goods of $\sigma_1$ hold four $D$-goods and three $S$-goods, so $9\in X$; and after the first eight goods of $\sigma_2$ the counts $(|D\cap P|,|S\cap P|,|X\cap P|)$ are $(3,3,2)$, so $11\in X$ as well, contradicting market balance on $Q_C$.
\end{proof}
Full enumeration confirms the counts (\appFifteen). By Theorem~\ref{thm:twelve} this size is minimal. Concatenating $h$ padded copies as contiguous modules forces agent~1 to hold two goods in at least $h$ of its triples, so no constant number of such shifts suffices for all~$m$.

\paragraph{All sizes under a one-singleton condition.}
Pad with bottom-ranked goods to a multiple of three. Call $\sigma_2$ \emph{one-singleton} if every prefix~$P$ of the padded $\sigma_2$ satisfies
\begin{equation}
 \#\{B\text{ a market triple}: |B\cap P|=1\}\le1.\label{eq:onesingleton}
\end{equation}
\begin{theorem}\label{thm:onesingleton}
Let $\sigma_0=\pi$ and let $\sigma_2$ be one-singleton. For every $\sigma_1$, the following rule gives in $O(m)$ time a market-block-balanced subjectively SD-EF1 allocation: in each market triple, agent~1 takes its $\sigma_1$-favourite good, agent~2 its $\sigma_2$-favourite remaining good, and agent~0 the last good.
\end{theorem}
\begin{proof}
In every triple agent~1's good precedes the other two in $\sigma_1$, so every $\sigma_1$ prefix gives agent~1 at least as many goods as each opponent. Agent~2 holds one of its two favourite goods of each triple. In a $\sigma_2$ prefix, a triple met in zero or three goods contributes nothing to an opponent's lead; a triple met in two goods contains agent~2's two favourites, one of which it holds; a triple met in one good contributes at most one. By~\eqref{eq:onesingleton} at most one triple is of the last kind. Market-block balance protects agent~0, and deleting the padding preserves every original prefix.
\end{proof}
For $3\mid m$ the condition characterizes exactly when this rule succeeds for every $\sigma_1$: if a prefix meets two triples in one good, put one chosen prefix good of every triple that it meets first in $\sigma_1$; agent~1 takes all of them and leads agent~2 by two or more. The class contains every ranking obtained from $\pi$ by disjoint adjacent transpositions, and, for $3\mid m$ and with the roles of agents~1 and~2 exchanged if necessary, every profile in which one real ranking lists each market triple contiguously (for instance $\sigma_1$ the reverse of $\pi$). The appendix extends this to short final blocks and to any number of agents: if all rankings but one are block-contiguous, a market-block-balanced subjectively SD-EF1 allocation exists (\appAllsizes). It also shows that if the goods split into consecutive market modules of 3, 6, 9 or 12 goods, each contiguous in both other rankings, then the modules can be solved separately with Theorem~\ref{thm:twelve} and the solutions concatenated in $O(m)$ time.

%%%%%%%%%%%%%%%%%%%%%%%%%%%%%%%%%%%%%%%%%%%%%%%%%%%%%%%%%%%%%%%%%%%%%%%%
\section{Barriers for the Matching Approach}\label{sec:barriers}
A perfect matching $S$ between market triples and $\sigma_1$ triples is easy to find, but it need not have a fair completion, and the results above suggest that completing or repairing an arbitrary matching cannot be the route to a general proof.

\begin{definition}[Prescribed completion]
Given $m=3q$ and rankings $\sigma_1,\sigma_2$ with $\sigma_0=\pi$, and a perfect matching $S$, decide whether a market-block-balanced subjectively SD-EF1 allocation with $A_1=S$ exists. The matching is a \emph{reset matching} if its good is last in every $\sigma_1$ triple, and satisfies the \emph{suffix restriction} if every good outside $S$ precedes every good of $S$ in $\sigma_2$.
\end{definition}

On $R=M\setminus S$ combine the residual market pairs and the residual $\sigma_1$ pairs; the union $H$ is an even alternating-cycle multigraph.
\begin{lemma}\label{lem:reset}
For a reset matching, the market-block-balanced fixed-$S$ allocations that satisfy agent~1's constraints are exactly the choices of one entire side of each component of $H$ for $X$, the other side going to~$D$.
\end{lemma}
\begin{proof}
Just before the $S$-good of the $k$-th $\sigma_1$ triple, $t=3k-1$ and $s=k-1$, so~\eqref{eq:one} forces $x=k$; successive equalities force exactly one $X$-good in each residual $\sigma_1$ pair, and market balance splits each residual market pair. Both edge systems are properly two-coloured, i.e.\ a side is chosen per component; conversely such choices satisfy all of agent~1's inequalities.
\end{proof}

\begin{theorem}\label{thm:hard}
Prescribed completion is NP-complete, even for reset matchings satisfying the suffix restriction.
\end{theorem}
\begin{proof}
Membership in NP is clear. We reduce from 3-SAT~\citep{karp1972}, repeating literals to pad short clauses. Let there be $k$ clauses. For a variable with $d$ occurrences create goods $a_j,b_j,z_j$ ($0\le j<d$) with $\sigma_1$ triples $(a_j,b_j,z_j)$ and market triples $(b_j,a_{j+1\bmod d},z_j)$, and put every $z_j$ into $S$. The residual graph is one cycle, and by Lemma~\ref{lem:reset} giving its $a$-side to $X$ means \emph{true}, its $b$-side \emph{false}. Add an \emph{anchor} cycle of the same form with $k+2$ positions and residual sides $\alpha_0,\dots,\alpha_{k+1}$ and $\beta_0,\dots,\beta_{k+1}$. For clause~$j$ let $L_{j1},L_{j2},L_{j3}$ be its occurrence goods (an $a$-good for a positive literal, a $b$-good for a negative one) and $\bar L$ the opposite good of the same occurrence. The residual part of $\sigma_2$ is $\alpha_0\alpha_1$, then for $j=0,\dots,k-1$ the gadget
\begin{equation}
L_{j1}\,L_{j2}\,L_{j3}\,\beta_j\,\bar L_{j1}\,\bar L_{j2}\,\bar L_{j3}\,\alpha_{j+2},\label{eq:gadget}
\end{equation}
then $\beta_k\beta_{k+1}$; $S$ comes last. The market triples are listed in any fixed order, which defines $\pi$. This gives $m=12k+6$ goods and strict rankings.

Before the suffix only agent~2's condition towards $D$ is non-automatic: the walk $W=x-d$ must stay at least $-1$. The first two goods force the anchor's $\alpha$-side into $X$ and give $W=2$. A clause gadget starts and ends at $W=2$. If it has $t$ true literals, $W$ equals $2t-1$ after its three literal goods and $2t-2$ after $\beta_j$, so a false clause violates fairness; for $t\ge1$ the walk stays at least $-1$ in the first part, and in the complement part, even taking its $t$ negative steps first, it stays at least $t-2\ge-1$, ends at $1$, and $\alpha_{j+2}$ restores $2$. The last two $\beta$ goods lead to $0$, and throughout the suffix $x=d=q$ and $s\le q$, so agent~2 stays fair. Hence a completion exists iff the formula is satisfiable.
\end{proof}
The natural linear relaxation is uninformative: giving every residual good weight $\tfrac12$ is always feasible for reset matchings with the suffix restriction (\appComplexity).

\paragraph{Fixed-$S$ infeasibility is not non-existence.}
Take $m=12$, $\sigma_1=(10,0,2,\,1,3,5,\,4,6,8,\,7,9,11)$, $\sigma_2=(10,1,0,3,6,9,4,7,2,5,8,11)$ and $S=\{2,5,8,11\}$. The residual graph is one cycle with sides $U=\{10,1,4,7\}$ and $V=\{0,3,6,9\}$, and $\sigma_2$ reads $UUVVVVUU$ on it. Giving $V$ to agent~2 fails after two goods, giving $U$ after six, while the half-half fractional solution passes; so $S$ has no fair completion. Yet $S'=\{2,5,8,9\}$ with $X=U$ and $D=\{0,3,6,11\}$ is a fair allocation. Completing a given matching is therefore the wrong question for a proof of existence, and hardness of completion is not evidence of hardness of existence.

\paragraph{Repair.}
On the reduction family the failure disappears if the matching may change. We first state the slack needed.
\begin{lemma}[Parallel-edge exchange]\label{lem:exchange}
Let $(D,S,X)$ be market-block-balanced with own-bundle balance, agent~1 SD-EF1, agent~2 SD-EF1 towards $S$, and $d-x\le2$ at every $\sigma_2$ prefix. Let $a\in S$ and $b\in D$ share a market triple and a $\sigma_1$ triple with $b$ earlier than $a$ in both real rankings, and let $I$ be the $\sigma_2$ prefixes containing $b$ but not $a$. If every prefix with $d-x=2$ lies in $I$ and $s\le x$ throughout $I$, swapping $a,b$ gives a subjectively SD-EF1 allocation preserving both block conditions.
\end{lemma}
\begin{proof}
Both blocks keep their counts. In the affected $\sigma_1$ prefixes $s$ grows by one and $d$ falls by one, improving both inequalities of agent~1. In $I$ the new counts satisfy $s'=s+1\le x+1$ and $d'=d-1\le x+1$; outside $I$ agent~2's inequalities already hold. Agent~0 is protected by market-block balance.
\end{proof}
\begin{theorem}\label{thm:repair}
Every profile produced by the reduction of Theorem~\ref{thm:hard} admits a market-block-balanced subjectively SD-EF1 allocation with own-bundle balance, whether or not the formula is satisfiable; a rankings-only algorithm finds one in $O(m)$ time, replacing at most one member of the prescribed~$S$.
\end{theorem}
\begin{proof}
Orient the variable cycles arbitrarily and the anchor positively. The clause walk never falls below $-2$, and its first $-2$, if any, is $\beta_j$ of the first false clause. This good shares both blocks with the anchor's $S$-good $z_j$ and precedes it in both rankings, so Lemma~\ref{lem:exchange} repairs every bad prefix with one exchange. The formula need not be recovered: read $S$ from the last $q$ positions of $\sigma_2$, build the residual cycles, and give agent~2 the side containing its first good on that component, which is the anchor.
\end{proof}
A replacement removes one good from $S$ and adds one, so two goods change owners. Theorem~\ref{thm:repair} does not make fixed-matching completion easy; it solves another task by declining to respect the prescribed matching, and it is guaranteed only on this family. Repair cannot be cheap in general.
\begin{theorem}\label{thm:distance}
For every $r\ge2$ there is a profile with $m=6r$ goods and $\sigma_0=\pi$, whose bipartite graph of market triples and $\sigma_1$ triples is simple and connected, and a reset matching $S$ formed by agent~2's bottom $m/3$ goods, such that (i)~$S$ has no fair completion; (ii)~every market-block-balanced subjectively SD-EF1 allocation with own-bundle balance has $|S\setminus A_1|\ge r$, with equality attained; (iii)~without own-bundle balance some such allocation has $|S\setminus A_1|=1$.
\end{theorem}
\begin{proof}
Let $q=2r$ and index the market triples by $\mathbb Z_q$; market triple $i$ is $(a_i,b_i,z_i)$. Let $\sigma_1$ be the concatenation over $j\in\mathbb Z_q$ of the triples $(a_{j-1},b_{j-2},z_j)$, and let $\sigma_2$ list $a_0,a_1,b_0,b_1,b_2,b_3$, then $a_2,\dots,a_{q-1}$, then $b_4,\dots,b_{q-1}$, then $z_0,\dots,z_{q-1}$; let $S=\{z_i\}$. Market triple $i$ thus meets $\sigma_1$ triples $i$, $i+1$, $i+2$ through $z_i,a_i,b_i$, and $S$ is a reset matching.

(i) The residual graph of $S$ is one alternating cycle with sides $\{a_i\}$ and $\{b_i\}$. By Lemma~\ref{lem:reset} agent~2 receives all $a_i$ or all $b_i$. With all $b_i$, agent~0 holds $a_0,a_1$ and agent~2 nothing after two goods; with all $a_i$, after six goods agent~2 holds two and agent~0 four. Both violate SD-EF1.

(ii) A perfect matching $S^*$ induces the permutation $i\mapsto i+\delta_i$ of $\mathbb Z_q$, where $\delta_i\in\{0,1,2\}$ says which of $z_i,a_i,b_i$ it uses; $S^*=S$ iff all $\delta_i=0$. The elements of a nontrivial cycle have $\delta_i\in\{1,2\}$, each replacing a member of $S$, and the increments along a cycle sum to a multiple of $q$, so a nontrivial cycle has at least $q/2=r$ elements. Since (i) shows $S^*\ne S$ for every fair own-bundle-balanced allocation, $|S\setminus A_1|\ge r$. The allocation $X=\{a_i\}$, $A_1=\{b_i:i\text{ even}\}\cup\{z_i:i\text{ odd}\}$ attains equality; (iii) is attained by $\widehat S=(S\setminus\{z_{q-1}\})\cup\{a_{q-1}\}$ with $\widehat X=(\{a_i\}\setminus\{a_2,a_{q-1}\})\cup\{b_2,b_{q-1}\}$. That both are fair, and the brute-force confirmation of (i)--(iii) for $r\le4$, are in \appDistance.
\end{proof}
Concatenating $h$ copies of the $r=2$ instance, each contiguous in every ranking, gives minimum replacement exactly $h$ without own-bundle balance and $2h$ with it, because for such modules global fairness is equivalent to fairness within every module. Hence no constant repair radius exists, with or without own-bundle balance.

\paragraph{Disjoint defects.}
For three agents and market order $0,\dots,m-1$ let $a_i(t)=|A_i\cap\{0,\dots,t-1\}|$. An allocation is \emph{near-block-balanced} if its three counts differ by at most one whenever $3\nmid t$; at $t=3k$ they are then either all $k$ or a permutation of $(k-1,k,k+1)$, and we call $t$ \emph{defective} in the second case; $D(A)$ is the set of defective boundaries.
\begin{proposition}\label{prop:defect}
Two near-block-balanced allocations $A,B$ are such that, for every additive market valuation consistent with $\pi$, at least one is market EF1 iff $D(A)\cap D(B)=\varnothing$.
\end{proposition}
\begin{proof}
With $\Delta_t=v_{t-1}-v_t\ge0$ ($v_m=0$), the violation of the comparison $i\to j$ equals $\sum_t c_{ij}(t)\Delta_t$ with $c_{ij}(t)=a_j(t)-a_i(t)-\mathbf 1[a_j(t)>0]$. This coefficient is at most $0$ at non-boundary prefixes, equal to $-1$ at non-defective boundaries, and at most $+1$ at defective ones. If the defect sets are disjoint, any row of $A$ plus any row of $B$ has coefficients summing to at most zero everywhere, so both violations cannot be positive. Conversely, the valuation that is $1$ on the first $t$ goods and $0$ elsewhere breaks both at a common defective~$t$.
\end{proof}
Such pairs need not exist. On a twelve-good profile (given in \appDefect), enumerating all $3^{12}=531{,}441$ allocations shows $1{,}080$ subjectively SD-EF1 ones, $58$ of them near-block-balanced, with defect sets $\{3,6\}$, $\{3,9\}$, $\{6,9\}$ and $\{3,6,9\}$ only, no two disjoint; the profile still has a cover by two allocations, one not near-block-balanced. The phenomenon already occurs for two agents: with the corresponding notion (prefix counts differing by at most two, defects at even prefixes) the eleven-good profile with identity market order,
\begin{align*}
\sigma_1&=(0,3,4,7,2,6,1,9,8,10,5),\\
\sigma_2&=(7,3,4,6,0,2,8,9,10,5,1),
\end{align*}
has 81 subjectively fair allocations, 24 of them near-block-balanced, and no disjoint pair, yet two allocations cover all additive valuations. Selection rules for the pair also fail (ten-good obstruction, appendix). These are failures of a proof route, not counterexamples to Open Question~1.

%%%%%%%%%%%%%%%%%%%%%%%%%%%%%%%%%%%%%%%%%%%%%%%%%%%%%%%%%%%%%%%%%%%%%%%%
\section{Evidence and Open Questions}\label{sec:discussion}
\paragraph{Evidence for general three-agent instances.}
A counterexample needs a \emph{hard} profile: no allocation that is subjectively SD-EF1 and market-block-balanced. Fixing the market order and treating agents as interchangeable, exhaustive enumeration of all $21{,}350{,}046{,}480$ ranking profiles for $m=7$, and all smaller ones, finds no hard profile. For larger $m$, hard profiles were generated by hill-climbing on the number of market-block-balanced fair allocations (874 independent restarts) followed by guided mutations. Each was decided exactly without enumerating allocations: a mixed-integer master problem seeks a market valuation defeating a finite set of subjectively fair allocations, and an integer program over all subjectively fair allocations returns one that survives, until a counterexample or a refuting set is found. A positive control with market EFX in place of EF1 finds counterexamples, each confirmed by enumeration. No counterexample appears among $4{,}087$ hard profiles with $13\le m\le21$ and $102$ hard profiles with up to $36$ goods built by concatenation; refuting sets had $3$ to $12$ allocations. These decisions rely on solver optimality up to a tolerance of $10^{-6}$; $90$ refuting sets, including the hardest profile of every search worker and 33 of the concatenated profiles, were re-verified by a branch search that solves only linear programs. Moreover, on every one of $90$ hard profiles saved from these searches, with up to $36$ goods, an integer program found two near-block-balanced subjectively fair allocations with disjoint market defects, hence a cover of all additive valuations by two allocations (Proposition~\ref{prop:defect}); such pairs cannot be guaranteed, as shown above. The evidence is heuristic above $m=7$ and is no proof.

\paragraph{Topology.}
Let the \emph{fairness complex} have the subjectively SD-EF1 allocations as vertices, two being adjacent when every bundle loses and gains at most one good. On all $18$ hard three-agent profiles with $10\le m\le12$ its first Betti number over $\mathbb F_2$ is between $4$ and $406$, so it is not simply connected, and the topological transfer principle of \citet{dupre2026balanced}, which requires simple connectivity for three agents, does not apply to it. This is not evidence against Open Question~1, only against this particular route.

\paragraph{Open questions.}
\begin{enumerate}
\item Does Open Question~1 hold for three agents, and does the market-aligned conjecture hold for all~$m$?
\item \emph{Two-cover question.} For two (or three) agents and additive markets, do two subjectively SD-EF1 allocations always cover all valuations? For two agents this holds on all $600$ random hard profiles with $8\le m\le13$ that we tested, and an adversarial search found no exception (\appFamily); for three agents it holds on the $90$ hard profiles above. A proof would strengthen Open Question~1 to a valuation-independent statement.
\item Is the $\lceil m/2\rceil$ of Theorem~\ref{thm:two} optimal for every profile, or only for the sampled ones?
\item How hard is finding a fair allocation when the matching may be chosen freely? Theorem~\ref{thm:hard} only concerns a prescribed one.
\end{enumerate}

\paragraph{Conclusion.}
A controlled swap path, built from the rankings alone, settles the two-agent question for every monotone market valuation, with extensions to chores and mixed items and a lower bound on the size of any cover. For market-aligned three-agent profiles, complete coverage proves existence through fifteen goods and a linear-time rule proves it for all sizes when one ranking is one-singleton. Completing a prescribed matching is NP-complete and repairing it can be expensive, so a proof for three agents must choose the matching, or avoid matchings altogether.

\paragraph{Use of generative AI\@.}
Following the AAMAS policy we disclose the use of generative AI. Large language models (ChatGPT, Codex, Claude, and Gemini for an initial literature scan) proposed conjectures and drafts of proofs, wrote verification code, and helped edit the text; the tools, their versions and the archived prompts are listed in the supplementary material. Every theorem here is either proved in the paper or marked computer-assisted. The proofs were re-derived line by line in separate audits, every finite claim was re-run locally, and the key computations were re-implemented independently; no proof assistant was used, and every citation was checked against the publisher's or arXiv's metadata. Code, logs and data are in the supplementary material and at \url{https://anonymous.4open.science/r/price-and-prize-0625}.

\bibliographystyle{ACM-Reference-Format}
\bibliography{references}

\end{document}
